% 原书物理页=印刷页24--25上；2.3全节，止于2.4标题前.
\section{Killing 型. Casimir 算子}\label{sec:ssl-2-3}
\begin{definition}\label{def:ssl-2-3-1}
设 $\mathfrak g$ 是一个 $K$-Lie 代数. $\mathfrak g$ 上的 $K$-双线性型 $B$ 称为\newterm{ssl-invariant-bilinear-form}{不变的}{invariant}，如果对一切 $X, Y, Z \in \mathfrak g$ 都有 $B(\operatorname{ad}_{Z}(X), Y) + B(X, \operatorname{ad}_{Z}(Y)) = 0$.
\end{definition}
\begin{example*}
不变双线性型的典范例子是 $\mathfrak{gl}_{n}(K)$ 上的型
\[(X, Y) \longmapsto \operatorname{tr}(XY).\]
这给出了构造许多其他不变双线性型的办法：对每个表示 $r : \mathfrak g \longrightarrow \mathfrak{gl}(V)$，$\mathfrak g$ 上的双线性型 $B_{r}(X, Y) := \operatorname{tr}(r(X)r(Y))$ 是不变的.
\end{example*}
\textbf{性质.}\quad 若 $B$ 是 $\mathfrak g$ 上的不变双线性型，而 $\mathfrak h \subset \mathfrak g$ 是 $\mathfrak g$ 的理想，则它的正交 $\mathfrak h^{\perp} = \{X \in \mathfrak g,\ \forall Y \in \mathfrak h,\ B(X, Y) = 0\}$ 也是 $\mathfrak g$ 的理想. 特别地，$\mathfrak g^{\perp}$ 是 $\mathfrak g$ 的理想.

\textbf{Killing 型.}\quad $\mathfrak g$ 上的\newterm{ssl-killing-form}{Killing 型}{Killing form} $B_{\operatorname{ad}}$ 是按上述方式与表示 $r = \operatorname{ad}$ 相伴的不变双线性型. 于是 $B_{\operatorname{ad}}(X, Y) = \operatorname{tr}(\operatorname{ad}_{X}\operatorname{ad}_{Y})$.

\begin{lemma}\label{lemma:ssl-2-3-2}
若 $\mathfrak h$ 是 $\mathfrak g$ 的理想\footnote{注意：若 $\mathfrak h$ 仅仅是 $\mathfrak g$ 的一个 Lie 子代数，则此结论不成立.}，则 $\mathfrak h$ 的 Killing 型是 $\mathfrak g$ 的 Killing 型在 $\mathfrak h$ 上的限制.
\end{lemma}
\begin{proof}
对 $X, Y \in \mathfrak h$，有 $\operatorname{ad}_{X}\operatorname{ad}_{Y}(\mathfrak g) \subset \mathfrak h$. 因此 $\operatorname{tr}(\operatorname{ad}_{X}\operatorname{ad}_{Y}|_{\mathfrak g}) = \operatorname{tr}(\operatorname{ad}_{X}\operatorname{ad}_{Y}|_{\mathfrak h})$.
\end{proof}

\begin{numpara}[非退化不变双线性型]\label{np:ssl-2-3-3}
现在设给定了 $\mathfrak g$ 上的一个非退化双线性型 $B$. 由它可造出泛包络代数的一个元素 $C_{B} \in \mathcal{U}\mathfrak g$，称为 Casimir 算子，定义为
\[C_{B} = \sum_{i=1}^{n} X_{i} X_{i}^{*},\]
其中 $X_{1}, \cdots, X_{n}$ 是 $\mathfrak g$ 的一组基，而 $X_{1}^{*}, \cdots, X_{n}^{*}$ 是 $\mathfrak g$ 关于非退化型 $B$ 的\newterm{ssl-dual-basis}{对偶基}{dual basis}.
\begin{lemma*}
上述定义与基的选取无关.
\end{lemma*}
\begin{proof}
设 $Y_{1}, \cdots, Y_{n}$ 是 $\mathfrak g$ 的另一组基. 有 $Y_{i} = \sum_{k} B(Y_{i}, X_{k}^{*}) X_{k}$ 与 $Y_{i}^{*} = \sum_{l} B(Y_{i}^{*}, X_{l}) X_{l}^{*}$. 于是
\begin{align*}
\sum_{i} Y_{i} Y_{i}^{*} &= \sum_{k,l} \Bigl( \sum_{i} B(Y_{i}, X_{k}^{*}) B(Y_{i}^{*}, X_{l}) \Bigr) X_{k} X_{l}^{*}\\
&= \sum_{k,l} B\Bigl( \sum_{i} Y_{i} B(Y_{i}^{*}, X_{l}), X_{k}^{*} \Bigr) X_{k} X_{l}^{*} = \sum_{k,l} B(X_{l}, X_{k}^{*}) X_{k} X_{l}^{*} = C_{B}.
\end{align*}
\end{proof}
\end{numpara}
% 原书p25上部，§2.3在§2.4标题之前的完整收尾.
\begin{example*}
取 $\mathfrak g=\mathfrak{gl}_n(K)$，$B$ 为由 $(X,Y)\mapsto\operatorname{tr}(XY)$ 定义的标准双线性型. 它非退化，且标准基 $(E_{ij})_{i,j}$ 的对偶基由 $E_{ij}^*=E_{ji}$ 给出.
\end{example*}
\begin{example*}
设 $\mathfrak g=\mathfrak{sl}_2(K)$，$B=B_{\mathrm{std}}$ 为由 $(X,Y)\mapsto\operatorname{tr}(XY)$ 定义的标准双线性型. 当 $\operatorname{car}(K)\neq2$ 时，容易验证它非退化，并且基 $(E,H,F)$ 的对偶基是 $(F,\tfrac12H,E)$. 这就重新得到上文引入的算子 $C=EF+\tfrac12H^2+FE$.
\end{example*}
\begin{exercise*}
\begin{enumerate}[label=\roman*)]
\item 验证 $(X,Y)\mapsto\operatorname{tr}(XY)$ 在 $\mathfrak{gl}_n(K)$ 以及经典 Lie 代数（$\mathfrak{so}$、$\mathfrak{sl}$、$\mathfrak{sp}$）上非退化.\footnote{译注：此问未附特征限制，按原文照录. 例如，当 $\operatorname{char}(K)$ 整除 $n$ 时，$\mathfrak{sl}_n(K)$ 中的非零标量矩阵与所有无迹矩阵的迹配对为零，因此一般情形需加限制.}
\item （$K=\mathbb R$ 或 $\mathbb C$.）证明：对 $\mathfrak{gl}_n(\mathbb R)$ 的每个在转置下稳定的 Lie 子代数 $\mathfrak g$，双线性型 $(X,Y)\mapsto\operatorname{tr}(XY)$ 都非退化.
\end{enumerate}
\end{exercise*}
我们已经在一个练习中见过，$\mathfrak{sl}_2$ 的 Casimir 算子的妙处在于它属于 $\mathcal U\mathfrak{sl}_2$ 的中心. 更一般地：
\begin{proposition*}
算子 $C_B$ 属于 $\mathcal U\mathfrak g$ 的中心.
\end{proposition*}
\begin{proof}
只需证明它与每个 $X\in\mathfrak g$ 交换. 而有
\begin{align*}
XC_B-C_BX
&=\sum_i(XX_iX_i^*-X_iX_i^*X)\\
&=\sum_i([X,X_i]X_i^*-X_i[X_i^*,X])\\
&=\sum_i\sum_k\bigl(B([X,X_i],X_k^*)X_kX_i^*\\
&\hspace{5em}{}-B(X_k,[X_i^*,X])X_iX_k^*\bigr)\\
&=\sum_{i,k}\bigl(B([X,X_k],X_i^*)+B(X_k,[X,X_i^*])\bigr)X_iX_k^*=0,
\end{align*}
最后一步来自双线性型 $B$ 的不变性.
\end{proof}
与 $\mathfrak{sl}_2$ 的情形相同，一个有用的推论是：对 $\mathfrak g$ 的任意表示 $(V,r)$，元素 $C$ 都诱导一个与 $\mathfrak g$ 的作用交换的算子 $r(C)=\sum_ir(X_i)r(X_i^*)$. 特别地，若 $V$ 不可约且 $K$ 代数闭（例如 $K=\mathbb C$），则由 Schur 引理，$r(C)$ 是一个位似.
